% !TEX root = ../report.tex

\section{Conclusion}
\label{sec:conclusion}

This paper asked whether diversification becomes fragile exactly when investors need it, and answered it with three measures conditioned on a stress regime that two independent definitions---a Markov-switching model on returns and a VIX-plus-drawdown rule---agree on. Three findings are robust. First, in realized terms diversification does weaken in stress: pairwise correlations rise and every asset universe concentrates onto fewer independent return directions, most extremely for international equities, whose leading principal component absorbs roughly $89\%$ of variance in stress. Second, the Forbes--Rigobon adjustment shows that this increase is volatility-driven \emph{interdependence} rather than structural \emph{contagion}---no pair, under either adjustment method and in any universe, survives as a robust breakdown. Third, and of most practical relevance, the fragility is concentrated in equities: international-equity pairs supply almost all of the volatility-bias-sensitive increases, while the cross-asset diversifiers---chiefly sovereign duration and gold---are genuinely resilient on average across the pooled stress regime, de-correlating further in most stress episodes, although the descriptive episode split shows the stock--Treasury hedge weakening and turning mildly positive in the inflation-driven 2022--2023 stress (Appendix Table~\ref{tab:stress_episode_hedges}); gold's diversifying role is stable across episode types.

The implication is diagnostic, not prescriptive. Diversification should be evaluated under stress regimes rather than on full-sample averages that blend calm and crisis, because the conditional covariance structure is both more concentrated than the unconditional one and partly inflated by a volatility effect that a naive correlation overstates; and the exposures that retain diversifying power when correlations rise are different assets, not other equities.

\paragraph{Further research.} Three extensions would tighten the evidence. First, \emph{dynamics}: this paper restricts attention to static regime comparisons, so a natural extension is to trace how comovement evolves after a stress onset---using a longer sample or an externally identified volatility shock (for example a high-frequency or narrative instrument) to obtain the statistical power that a handful of ex-post regime transitions cannot provide. Second, \emph{dependence model}: the single-factor, variance-only Forbes--Rigobon correction could be generalized to a multi-factor or dynamic-conditional-correlation specification and complemented by tail-dependence or copula measures, which capture the asymmetric, extreme-state comovement that a Gaussian correlation misses. Third, \emph{economic value}: whether the regime-conditional, variance-adjusted dependence estimates translate into out-of-sample resilience---for instance a regime-aware minimum-variance or risk-parity allocation benchmarked against an unconditional one---is the natural test of whether these diagnostics are useful to an investor rather than only descriptive.
